%=====================================================================
\clearpage
\section{Erweiterung ab Version 2.37: NURBS-Profil in 1D-Gittern (M24)}\label{sec:nurbs1d}
%=====================================================================
Ein Liniengitter mit Überhang oder Hinterschnitt lässt sich heute nur über eine STEP-Datei beschreiben
(\verb|cad|-Schicht, M23); \verb|profile1d| kennt nur einwertige Profile $y=g(x)$. M24 ergänzt \verb|profile1d| um
\verb|"shape": "nurbs"|: \emph{eine} durchgehende, periodische NURBS-Linie in der $(x,y)$-Ebene trennt das Material
\verb|below| (unterhalb) vom Material \verb|above|. Die Linie darf Überhänge haben, sich aber nicht selbst schneiden.
Alles Weitere sind Bausteine von M23 ohne Nurbforge: Zerlegung in rationale Bézierstücke (Piegl--Tiller A5.6,
\verb|Contour2D::nurbs|), Linienschnitt (\verb|lineCrossings|), Normalenfeld \eqref{eq:cad1d_N} und
\verb|Sectioned1DLayer| (RK, RCWA-Treppe, Tschebyschow, ASR, Nahfeld, Gradienten). Anlass: Wunsch aus der
Anwendung (Issue \#30).

\subsection{Annahmen und Festlegungen}
\begin{enumerate}[nosep]
  \item Neuer Wert \verb|"shape": "nurbs"| des bestehenden Typs \verb|profile1d|; die übrigen Formen bleiben unverändert
        und bitgleich. Nur in 1D-Stapeln.
  \item \textbf{Eine Schicht.} Die Schicht bleibt nach außen \emph{eine} Schicht mit internen Bruchstellen in $y$
        (Abschnitt~\ref{sec:nurbs1d_bands}); JSON-Pointer späterer Schichten, Sweeps und Berichte verschieben sich
        nicht. Eine Expansion in mehrere Schichten wie bei \verb|cad| entfällt.
  \item Die Bänder zwischen den Bruchstellen sind Integrationsintervalle, keine Treppe: RK wertet $Q(y)$ an jedem
        Stufenknoten aus dem exakten Linienschnitt aus. Eine Treppe entsteht nur mit RCWA (\verb|slices|).
  \item Faktorisierung wie \verb|Sectioned1DLayer| (Schlüssel \verb|factorization| der Schicht wie bei allen
        Schichttypen): isotrop \verb|laurent|, \verb|FFF_iso|, \verb|FFF_sym| (Standard),
        anisotrop \verb|laurent|, \verb|FFF_B| (Standard). \verb|FFF_rig| und \verb|FFF_A| verlangen ein
        vorzeichenbehaftetes Normalenfeld und werden abgelehnt, auch wenn die Kurve ein Graph ist (für Graphen
        bleibt \verb|profile1d| \verb|points|/\verb|Function|).
  \item Keine Formgradienten (nach \verb|depth_um| oder Kontrollpunkten) und keine Funktionen auf Vorrat. Sweeps
        über \verb|Simulation::setParameter| bauen die Schicht aus dem geänderten JSON neu und sind daher korrekt.
\end{enumerate}

\subsection{Konfiguration}
{\footnotesize\begin{verbatim}
{"type": "profile1d", "shape": "nurbs", "depth_um": 0.25, "below": "si", "above": "air",
 "nurbs": {"degree": 3, "knots": [0, 0, 0, 0, 0.3, 0.6, 1, 1, 1, 1],
           "points_um": [[0.0, 0.05], [0.1, 0.25], [0.2, 0.2], [0.3, 0.0], [0.4, 0.1], [0.5, 0.05]],
           "weights": [1, 1, 0.8, 1, 1, 1]},
 "slices": 0, "chebyshev_nodes": 0}
\end{verbatim}}
\begin{itemize}[nosep]
  \item \verb|points_um|: Kontrollpunkte $(x,y)$ in \textmu m, $y$ relativ zur Unterkante der Schicht. Knotenvektor
        geklemmt (erster und letzter Knoten mit Vielfachheit $p+1$), nicht fallend, Länge $n+p+2$ bei $n+1$ Punkten.
        \verb|weights| optional ($>0$, Standard 1). Grad $1\le p\le9$ (Grenze von \verb|BezierSegment|).
  \item \textbf{Periodizität:} Die Kurve ist offen und läuft über genau eine Periode. Wegen des geklemmten
        Knotenvektors sind ihre Enden der erste und der letzte Kontrollpunkt; verlangt ist
        $\mathbf P_{\text{Ende}}=\mathbf P_{\text{Anfang}}+(d_x,0)$ mit der Toleranz von \verb|Contour2D::validate|
        ($10^{-12}$ der Ausdehnung). Ein Höhenunterschied ist ein Fehler (es wird keine Wand eingefügt), ebenso ein
        Versatz $\ne d_x$ in $x$ (auch $-d_x$: die Kurve muss netto in $+x$ laufen). Der Anfang darf bei beliebigem
        $x$ liegen; zwischendurch darf die Kurve links von $0$ oder rechts von $d_x$ verlaufen.
  \item \textbf{Naht:} Die Tangenten an Anfang und Ende müssen nicht übereinstimmen; ein Knick an der Naht ist eine
        Ecke wie ein C$^0$-Knoten im Inneren (kritische Höhe).
  \item \verb|depth_um| (Dicke der Schicht) ist optional. Ohne Angabe gilt $\verb|depth_um|=\max_t y(t)$ (exakt
        aus den kritischen Höhen, nicht aus den Kontrollpunkten), wie bei den übrigen \verb|profile1d|-Formen, deren
        Profil die Schicht genau ausfüllt. Mit Angabe darf über der Linie eine Schicht aus \verb|above| bleiben.
        In beiden Fällen muss $0\le y(t)\le\verb|depth_um|$ für alle $t$ gelten (exakt über die $y$-Extrema geprüft);
        ein Bereich $0\le y<\min y$ ist zulässig und besteht aus \verb|below|.
  \item \verb|slices| (RCWA-Treppe, 0: Löserstandard) und \verb|chebyshev_nodes| (Tschebyschow-Interpolation in $y$ je
        Band, 0: aus) wie bei \verb|Sectioned1DLayer|.
\end{itemize}

\subsection{Linienschnitt einer offenen periodischen Kurve}\label{sec:nurbs1d_line}
Die Kurve wird nach A5.6 in rationale Bézierstücke $\mathbf C_s(t)$, $t\in[0,1]$, $s=1,\dots,S$ zerlegt. Der neue
Provider \verb|NurbsProfileLineProvider| (\verb|LineSectionProvider|) schneidet bei der Höhe $y=c$ (nie auf einer
kritischen Höhe):
\begin{enumerate}[nosep]
  \item Schnittpunkte je Stück mit \verb|lineCrossings| auf $t\in[0,1)$. Das Ende des letzten Stücks ist der um $d_x$
        verschobene Anfang des ersten; ein Schnittpunkt genau an der Naht zählt daher genau einmal.
  \item Normale \eqref{eq:cad1d_n} mit Vorzeichen, $\mathbf n=(\dot y,-\dot x)/|\dot{\mathbf C}|$. Die Kurve läuft
        netto in $+x$; \verb|below| liegt rechts von der Laufrichtung (im Uhrzeigersinn), denn sie begrenzt das
        Gebiet unter sich. An einem Schnittpunkt zeigt die rechte Seite in Richtung $\mathbf n$, also ist das Material
        unmittelbar \emph{rechts} des Schnittpunkts $x_k$
        \begin{equation}
          \text{mat}(x_k^+)=\begin{cases}\verb|below|, & \dot y(t_k)>0\ (n_x>0),\\ \verb|above|, & \dot y(t_k)<0\ (n_x<0).\end{cases}
          \label{eq:nurbs1d_rule}
        \end{equation}
  \item Rückfaltung $x\mapsto x-d_x\lfloor x/d_x\rfloor$ in $[0,d_x)$, sortieren. \textbf{Konsistenz:} Die Materialien
        rechts aufeinanderfolgender Schnittpunkte müssen sich abwechseln (zyklisch, also gerade Anzahl). Sonst ist die
        Kurve selbstschneidend (auch mit ihrem periodischen Bild) $\to$ Fehler \verb|NURBS_SELF_INTERSECTION|
        mit der Höhe in \textmu m.
  \item Schnittpunkte näher als $10^{-10}d_x$ (nach der Rückfaltung): gleiche Richtung ist derselbe Schnittpunkt
        zweimal gefunden (Naht: $t=0$ des ersten und $t\to1$ des letzten Stücks), entgegengesetzte Richtung begrenzt
        einen Splitter schmaler als die Toleranz neben einem Extremum (z.\,B.\ ein Tal auf der Unterkante, $10^{-14}$
        darüber abgetastet); er entfällt.
  \item Ohne Schnittpunkt liegt $c$ unter oder über der ganzen Kurve: \verb|below| für $c<\min y$, sonst \verb|above|
        (nach dem Wegfall eines Splitters: \verb|below| in der unteren Hälfte von $[\min y,\max y]$).
  \item Ergebnis wie \verb|lineSection| (Abschnitt~\ref{sec:cad1d_line}): Bruchstellen $0=\xi_0<\dots<\xi_K=d_x$, ein
        Material je Intervall, Normalen an den Materialgrenzen. Eine Grenze bei $0\equiv d_x$ bleibt als
        $\xi_0$ mit ihrer Normalen stehen; gleiche Nachbarn verschmelzen (nur über die Naht möglich).
\end{enumerate}
Ein geschlossener \verb|ProfileBody| (Kurve plus senkrechte Kanten bei $x=0$, $x=d_x$ und Boden) wäre nur ohne Überhang
über die Zellnaht korrekt; der eigene Provider braucht diese Einschränkung nicht. Der Vergleich mit einem solchen
Körper dient als Test (V25a).

\subsection{Kritische Höhen und Bänder}\label{sec:nurbs1d_bands}
$Q(y)$ ist nicht glatt, wo Schnittpunkte entstehen oder verschwinden (wurzelförmig) oder die Normale springt. Kritische
Höhen sind daher
\begin{enumerate}[nosep]
  \item die $y$-Extrema im Inneren jedes Stücks: Mit $y_s(t)=N(t)/D(t)$, $N=\sum_jw_jy_jB^p_j$, $D=\sum_jw_jB^p_j$, ist
        \begin{equation}
          \dot y_s=\frac{N'D-ND'}{D^2},\qquad N'D-ND'=\sum_{k=0}^{2p-1}c_kB^{2p-1}_k,\quad
          c_k=\sum_{i+j=k}\frac{\binom{p-1}{i}\binom{p}{j}}{\binom{2p-1}{k}}\,p\bigl(\Delta N_iD_j-\Delta D_iN_j\bigr)
          \label{eq:nurbs1d_crit}
        \end{equation}
        mit $N_j=w_jy_j$, $D_j=w_j$, $\Delta N_i=N_{i+1}-N_i$ (Produkt zweier Bernstein-Polynome). Nullstellen in
        $(0,1)$ über Bernstein-Unterteilung und abgesicherten Newton wie \verb|lineCrossings|; Rückgabe $y_s(t)$.
  \item die Höhen aller Stückgrenzen (Knoten im Inneren, auch C$^1$/C$^2$-Knoten, und die Naht). Das schließt Knicke
        (C$^0$-Knoten), Extrema an Knoten und waagrechte Stücke ein; überzählige Grenzen kosten nur einen RK-Schritt.
\end{enumerate}
Werte innerhalb $10^{-12}\,\verb|depth_um|$ werden zusammengelegt; Werte $\le0$ und $\ge\verb|depth_um|$ entfallen.
Ergebnis: $0<\eta_1<\dots<\eta_B<\verb|depth_um|$, die internen Bruchstellen der Schicht.
\begin{itemize}[nosep]
  \item \verb|Sectioned1DLayer| erhält \verb|yBreaks_um| (leer für M23). \verb|StackLayer::yBreaks_um()| (virtuell,
        Standard leer) gibt sie an den Löser.
  \item RK (Löser und Gradient, Plan eq:slicecrit unverändert): die Stabilitätsscheiben werden zusätzlich an den
        Bruchstellen geteilt, wie an den Grenzen der Hintergrundschichten (\verb|bgCuts|). RK wertet $Q$ auch an den
        Enden eines Schrittes aus; an einer Bruchstelle kann $Q$ springen (waagrechtes Stück) oder der Schnitt
        tangential sein. Daher wird $Q$ in jeder geschnittenen Scheibe $[y_a,y_b]$ bei
        $\operatorname{clamp}(y,\,y_a+\delta,\,y_b-\delta)$, $\delta=10^{-9}(y_b-y_a)$, ausgewertet: der einseitige
        Grenzwert; bei stetigem $Q$ heben sich die Fehler der beiden Nachbarschritte in erster Ordnung auf. Der
        Gradient verwendet dieselbe Klemmung an allen Stufenknoten (vorwärts und rückwärts). Ohne Bruchstellen ändert
        sich nichts (bitgleich).
  \item Ein Schnitt genau auf einer Bruchstelle (RCWA-Scheibenmitte, Nahfeldhöhe) wird $10^{-9}\,\verb|depth_um|$
        darüber genommen.
  \item Tschebyschow: ein Interpolant je Band (Knoten und Fehlerkontrolle je Band wie M14), nie über eine Bruchstelle.
  \item RCWA: die Treppe bleibt gleichmäßig über die ganze Schicht (Mittelpunktsregel; die Bruchstellen ändern ihre
        Konvergenz nicht).
  \item Eine Mindestdicke der Bänder (\verb|min_band_um|, \verb|FEATURE_SMALL|) wie im Slicer entfällt: Ein dünnes
        Band ist ein kurzer RK-Schritt, keine eigene Schicht. Bänder unter $\min y$ und über $\max y$ sind homogen;
        sie werden nicht als eigene Schichtart erkannt, der Schnitt liefert dort ein Material.
\end{itemize}

\subsection{Normalen und Faktorisierung}
Die Normalen an den Schnittpunkten sind exakt \eqref{eq:cad1d_n} (Newton auf Maschinengenauigkeit in
\verb|lineCrossings|), dazwischen gilt das Feld \eqref{eq:cad1d_N} mit den Koeffizienten in geschlossener Form
(\verb|normalField|). Mit ASR in $x$ (M17) werden die Bruchstellen wie in M23 über \verb|mapIntervals| abgebildet und
die Normale an der physikalischen Stelle ausgewertet; der Weg ist derselbe wie bei \verb|cad| in 1D.

\subsection{Prüfungen}
Alle beim Laden, mit JSON-Pointer der Schicht und, wo sinnvoll, Ort in \textmu m; Kennungen in eckigen Klammern am
Anfang der Meldung:
\begin{center}
\small
\begin{tabular}{p{4.6cm}p{9.8cm}}
\toprule
Kennung & Fall \\
\midrule
\verb|NURBS_INPUT| & Grad außerhalb 1..9, zu wenige Punkte, Knotenzahl $\ne$ Punkte $+$ Grad $+1$, Knoten fallend oder
                     nicht geklemmt, Gewicht $\le0$, Punkt nicht $[x,y]$ \\
\verb|NURBS_NOT_PERIODIC| & $y$ an Anfang und Ende verschieden \\
\verb|NURBS_PERIOD| & Versatz in $x$ nicht $+d_x$ (auch $-d_x$: falsche Orientierung) \\
\verb|NURBS_RANGE| & $\min y<0$ oder $\max y>\verb|depth_um|$, \verb|depth_um| $\le0$ \\
\verb|NURBS_SELF_INTERSECTION| & echter Schnitt der abgeflachten Kurve mit sich oder ihren Bildern bei $\pm d_x$
                     (\verb|properCrossing|, Toleranz $10^{-9}d_x$), oder Materialien wechseln nicht ab (in jeder
                     Bandmitte geprüft) \\
\verb|NURBS_STACK| & 2D-Stapel, Zylinderstapel, PML (\verb|aperiodic|) \\
\bottomrule
\end{tabular}
\end{center}
Die Faktorisierung prüft das Modell (Meldung von \verb|Sectioned1DLayer|). Gradienten nach \verb|depth_um| oder
\verb|nurbs/...| lehnt \verb|gradient| mit den bestehenden Meldungen ab.

\subsection{Tests und Abnahme (M24)}
\begin{itemize}[nosep]
  \item \textbf{V25a} Unit (\verb|tests/unit/test_nurbs_profile.cpp|): Die Bézierstücke reproduzieren die NURBS-Kurve
        (de Boor) auf $10^{-14}$. Kritische Höhen \eqref{eq:nurbs1d_crit} für eine kubische Polynomkurve und einen
        Kreisbogen (rationale Quadrik) gegen die geschlossene Form ($10^{-12}$). Provider gegen \verb|lineSection| eines
        gleichwertigen geschlossenen \verb|ProfileBody| an vielen Höhen (Bruchstellen und Normalen $10^{-13}$,
        Materialien gleich), auch mit Überhang und mit Kurvenstart bei beliebigem $x$. Alle Fehlerfälle der Tabelle.
  \item \textbf{V25b} Grad 1 durch die Punkte eines \verb|profile1d| \verb|points| bzw.\ \verb|triangle|: mit Laurent
        $Q$ an fünf Höhen gleich ($10^{-13}$), Effizienzen RCWA $10^{-12}$, RK $10^{-10}$; isotrop und anisotrop.
  \item \textbf{V25c} Gerundetes Trapez aus V24c ($CD=100$\,nm, $h=100$\,nm, $15^\circ$-Wände, $r=10$\,nm) exakt als
        NURBS mit Kreisbögen gegen den \verb|Function|-Wert von V24c (RK, Laurent): $10^{-10}$, angestrebt das Niveau
        $2{,}6\cdot10^{-13}$. Im CI-Job \verb|cad| zusätzlich gegen \verb|examples/cad_line_grating_1d.json| (RK, Toleranz $10^{-10}$;
        verschiedene Schrittfolgen, eine Schicht gegen drei: $10^{-9}$, gemessen $3{,}2\cdot10^{-10}$).
  \item \textbf{V25d} Überhang (Pilzprofil): $R+T=1$ ($10^{-10}$, verlustfrei), Reziprozität (fällt mit $N_x$),
        RCWA $\to$ RK mit steigender \verb|slices|-Zahl, Konvergenz in $N_x$. Kurvenstart verschoben (dieselbe
        Geometrie ab einem inneren $x$): Effizienzen $10^{-12}$.
  \item \textbf{V25e} Naht: glatte Kurve und Kurve mit Knick an der Naht, je ab $x=0$ und ab innerem $x$ beschrieben:
        $10^{-12}$. Ein Schnitt genau bei $y_0$ liefert die Grenze bei $0\equiv d_x$ genau einmal.
  \item \textbf{V25f} Tschebyschow an/aus ($10^{-8}$), ASR in $x$ (gegen ohne, Konvergenz), Nahfeld läuft; ohne
        Bruchstellen ändert sich nichts (alle Golden Files bitgleich).
  \item \textbf{V25g} Fehler: Kennungen der Tabelle mit Pointer; Gradient nach \verb|depth_um| und Kontrollpunkt
        abgelehnt; Sweep über einen Kontrollpunkt gleich der Neuberechnung.
  \item \textbf{V25h} Beispiel \verb|examples/profile1d_nurbs.json| mit eigenem Golden File; Python-Test (JSON läuft
        durch, Ergebnis gleich \verb|fgsolve|).
\end{itemize}

\subsection{Umsetzung (v2.37)}
\begin{itemize}[nosep]
  \item \verb|src/geometry/nurbs_profile.cpp|: \verb|NurbsProfile| (Zerlegung über \verb|Contour2D::nurbs|, Prüfungen,
        Extrema nach \eqref{eq:nurbs1d_crit} mit eigener Bernstein-Nullstellensuche, da $2p-1$ den Grad 9 von
        \verb|BezierSegment| übersteigen kann), \verb|NurbsProfileLineProvider|. Konfiguration in
        \verb|parseNurbsProfile| (die Schicht braucht $d_x$ und die Stapelart, daher im Schichtlauf wie \verb|cad|).
  \item \textbf{Befund M23:} \verb|lineCrossings| verlor eine innere Nullstelle, wenn die Kurve die Gerade genau am
        Ende eines Stücks erreicht: Der Endkoeffizient 0 wurde durch den Nachbarn ersetzt, \verb|refine| sah am Ende
        aber $y=c$ und gab das Ende zurück. Jetzt erhält \verb|refine| die Vorzeichen der ersetzten Enden. In M23
        trat der Fall nicht auf (Schnitte nie genau auf Stückgrenzen), in M24 beim Schnitt genau durch die Naht.
  \item Ergebnisse (\verb|docs/validation.md|): V25b Grad 1 gegen \verb|points|/\verb|triangle| RCWA $\le1{,}4\cdot10^{-15}$,
        RK $\le3{,}1\cdot10^{-13}$; V25c gerundete Linie gegen \verb|Function| $2{,}1\cdot10^{-13}$ (wie V24c mit
        STEP); V25d Pilz $R+T=1$ auf $10^{-14}$, RCWA $\to$ RK $2{,}5\cdot10^{-3}$, $1{,}4\cdot10^{-5}$,
        $2{,}3\cdot10^{-6}$ bei 20, 80, 320 Scheiben, Anfang verschoben $10^{-15}$; V25e Naht $\le1{,}3\cdot10^{-14}$;
        V25f Tschebyschow $4{,}8\cdot10^{-15}$.
  \item Abnahme angepasst: Mit ASR ist $R+T=1$ nur bis auf die Umrechnung der Ordnungen erfüllt ($\sim N^{-4}$, wie
        M17; $2\cdot10^{-7}$ bei $N=16$). Die Konvergenz in $N_x$ ist mit \verb|FFF_sym| nicht monoton (schon bei
        $N=8$ auf $2\cdot10^{-5}$ am Wert für $N=56$); geprüft wird daher, dass \verb|FFF_sym| bei $N=8$ zehnmal näher
        liegt als Laurent und Laurent monoton folgt ($6{,}4\cdot10^{-4}$, $2{,}3\cdot10^{-4}$ bei $N=8$, 24).
\end{itemize}

\subsection{Reihenfolge der Umsetzung}
\begin{enumerate}[nosep]
  \item Plan, dann Tests V25a (Unit).
  \item \verb|src/geometry/nurbs_profile.cpp|: Zerlegung, kritische Höhen \eqref{eq:nurbs1d_crit}, Prüfungen,
        \verb|NurbsProfileLineProvider| mit Regel \eqref{eq:nurbs1d_rule}.
  \item \verb|Sectioned1DLayer::yBreaks_um|, \verb|StackLayer::yBreaks_um()|, Schnitte in RK (Löser und Gradient),
        Tschebyschow je Band (Golden Files bitgleich).
  \item Konfiguration (\verb|profile1d| \verb|nurbs| $\to$ \verb|Sectioned1DLayer|), Prüfungen mit Pointer; V25b--h.
  \item Python (\verb|LLM_REFERENCE.md|), Beispiel, Handbuch, Dokumentation.
\end{enumerate}
